\newpage
\phantomsection
\label{prf:induction-well-ordering-equivalent}
\LRAProofFor{prop:induction-well-ordering-equivalent}

\begin{remark*}[Return]
\hyperref[prop:induction-well-ordering-equivalent]{Return to Proposition}
\end{remark*}

\begin{proposition*}[Induction and Well-Ordering Are Equivalent]
Over the axioms of $\mathbb{N}$ without P5, the weak induction principle,
the well-ordering principle, and the strong induction principle are
logically equivalent.
\end{proposition*}

\begin{proof}[Professional Standard Proof]
\LRAProofBodyStart
It suffices to prove weak induction implies well-ordering and
well-ordering implies weak induction; the equivalence with strong
induction follows from the standard mutual derivations of weak and strong
induction.

Assume weak induction and suppose, toward contradiction, that a nonempty
set $S \subseteq \mathbb{N}$ has no least element. Let
$P(n)$ be the assertion $n \notin S$. Since $S$ has no least element,
$0 \notin S$, so $P(0)$ holds. If $P(k)$ holds for all $k \le n$, then
$n+1 \notin S$, because otherwise $n+1$ would be the least element of $S$.
Thus induction gives $P(n)$ for every $n$, contradicting that $S$ is
nonempty. Hence every nonempty subset of $\mathbb{N}$ has a least element.

Conversely, assume well-ordering. Suppose $P(n_0)$ holds and
$P(n) \Rightarrow P(n+1)$ for every $n \ge n_0$. If some $P(n)$ fails, then
$S=\{n \ge n_0 : \neg P(n)\}$ is nonempty and has a least element $m$.
Since $P(n_0)$ holds, $m>n_0$, so $m-1 \ge n_0$. By minimality of $m$,
$P(m-1)$ holds, and the inductive step gives $P(m)$, a contradiction.
Therefore $P(n)$ holds for all $n \ge n_0$.
\end{proof}

\begin{proof}[Detailed Learning Proof]
\LRAProofBodyStart
The two principles encode the same obstruction. Induction says there is no
first place where a hereditary property can fail after its base case.
Well-ordering says that any nonempty set of failures would have a first
failure. Turning one sentence into the other gives the equivalence.
\end{proof}

\begin{remark*}[Proof structure]
Show each principle rules out the same minimal counterexample pattern.
\end{remark*}

\begin{dependencies}
\begin{itemize}
  \item Natural-number order.
  \item Weak induction principle.
  \item Strong induction principle.
  \item Well-ordering principle.
\end{itemize}
\end{dependencies}

\clearpage
